\section[Illustrations \& References]{Illustrations \& References}

\blockframe{4}{Illustrations \& References}{Tables, figures and citations}

\begin{frame}
    \frametitle{Where we are going}
    \vfill
    \begin{roadmap}{4}
        \roadmapresults
    \end{roadmap}
    \vfill
\end{frame}

\begin{frame}
    \frametitle{Block 4 in one line}
    \vfill
    \begin{center}
        {\LARGE No number is typed by hand\par}
        \vspace{1cm}
        \uncover<2->{%
            {\large\color{black!60}
            Sections \textperiodcentered{} equations \textperiodcentered{} tables
            \textperiodcentered{} figures \textperiodcentered{} citations\par}
        }
        \bigskip
        \uncover<3->{%
            {\large\texttt{\textbackslash label\{name\}} \quad then \quad \texttt{\textbackslash ref\{name\}}\par}
        }
    \end{center}
    \vfill
    \note{%
        Say it: ``you never type a number that \LaTeX{} can work out for you.''
        Then stop talking for a beat.
        You give the thing a label, and you ask for it by name.
        \begin{itemize}
            \item The next twenty minutes are one idea wearing four hats.
                If they leave with only this, block 4 worked.
            \item They have seen it twice already: the table of contents in block 1 and \texttt{\textbackslash eqref} in block 3.
                This names a pattern they have felt.
                Point that out; it is why it will stick.
        \end{itemize}
    }
\end{frame}

% The highlighted lines follow the bullet points.
\begin{frame}
    \frametitle{A table, the hard way first}
    \only<1,5-| handout:1,4->{\tcbinputlst{listings/04-tabular.tex}}%
    \only<2| handout:2>{\tcbinputlst[highlightlines={1}]{listings/04-tabular.tex}}%
    \only<3-4| handout:3>{\tcbinputlst[highlightlines={2-4}]{listings/04-tabular.tex}}%
    \begin{columns}[T, onlytextwidth]
        \begin{column}{.52\textwidth}
            \begin{outputbox}
                \footnotesize
                \begin{tabular}{lllr}
    Vessel & Material & Lid & Time constant \\
    Mug & glazed ceramic & no & 24 min \\
    Glass & borosilicate & no & 18 min \\
    Flask & steel & yes & 290 min \\
\end{tabular}

            \end{outputbox}
        \end{column}
        \begin{column}{.44\textwidth}
            \begin{itemize}
                \item<2-| handout:2-> \texttt{\{lllr\}}: one letter per column
                    (\textbf{l}eft, \textbf{c}entre, \textbf{r}ight)
                \item<3-| handout:3-> \texttt{\&}: next column
                \item<4-| handout:3-> \texttt{\textbackslash\textbackslash}: end of row
            \end{itemize}
        \end{column}
    \end{columns}
    \vfill
    \begin{alertblock}<5-| handout:4>{}
        Not a table yet: no caption, no number, no placement
    \end{alertblock}
    \note{%
        \begin{itemize}
            \item \texttt{\textbackslash\textbackslash} here is its other legitimate home (block 2 promised two).
            \item A bare \texttt{tabular} lands wherever it falls, mid-sentence or across a page break.
            \item Let them see the bare \texttt{tabular} misbehave.
                The TODO has them type it unwrapped on purpose,
                so the float on the next slide solves a problem they actually just had.
            \item \texttt{p\{3cm\}} for wrapping columns is in the handout.
                Mention it when someone's text runs off the page, not before.
        \end{itemize}
    }
\end{frame}

\begin{frame}
    \frametitle{Make it float}
    \begin{columns}[T, onlytextwidth]
        \begin{column}{.52\textwidth}
            \only<1,5-| handout:1,2>{\tcbinputlst[fontsize=\footnotesize]{listings/04-table-float.tex}}%
            \only<2| handout:none>{\tcbinputlst[fontsize=\footnotesize, highlightlines={1,8}]{listings/04-table-float.tex}}%
            \only<3| handout:none>{\tcbinputlst[fontsize=\footnotesize, highlightlines={3}]{listings/04-table-float.tex}}%
            \only<4| handout:none>{\tcbinputlst[fontsize=\footnotesize, highlightlines={4}]{listings/04-table-float.tex}}%
        \end{column}
        \begin{column}{.44\textwidth}
            \begin{itemize}
                \item<2-| handout:2> \texttt{[htbp]}: \textbf{h}ere, \textbf{t}op, \textbf{b}ottom, own \textbf{p}age
                \item<3-| handout:2> \texttt{\textbackslash caption}: numbered for you
                \item<4-| handout:2> \texttt{\textbackslash label} \textbf{after} \texttt{\textbackslash caption}
                \item<5-| handout:2> \texttt{\textbackslash usepackage\{caption\}}: proper spacing
            \end{itemize}
        \end{column}
    \end{columns}
    \vfill
    \begin{exampleblock}<6-| handout:2>{}
        Let \LaTeX{} decide where floats go
    \end{exampleblock}
    \note{%
        \begin{itemize}
            \item \texttt{[htbp]} is a wish list, tried in order.
            \item Fighting float placement is a huge time sink; do not start.
            \item Without the \texttt{caption} package the gap above a table is far too tight.
            \item \texttt{\textbackslash label} before \texttt{\textbackslash caption} labels the wrong thing.
                It is the rule that will bite,
                because it fails silently: a reference to the wrong number, no error.
                Say it twice.
            \item \texttt{[H]} from the \texttt{float} package: say it exists, say it is a last resort, move on.
                Taught as normal, it gives documents full of half-empty pages.
            \item The 2025 deck said \texttt{!} ``forces'' placement.
                It does not, it only relaxes \LaTeX{}'s internal limits.
                \texttt{[H]} is the one that forces.
        \end{itemize}
    }
\end{frame}

% Both tables are typeset for real from the same files as the code.
\begin{frame}
    \frametitle{Lines in tables}
    \begin{columns}[T, onlytextwidth]
        \begin{column}{.5\textwidth}
            \tcbinputlst[fontsize=\footnotesize]{listings/04-grid.tex}
        \end{column}
        \begin{column}{.45\textwidth}
            \begin{center}
                \begin{tabular}{|l|l|l|}
    \hline
    Vessel & Material       & Lid \\
    \hline
    Mug    & glazed ceramic & no  \\
    \hline
    Flask  & steel          & yes \\
    \hline
\end{tabular}

            \end{center}
            \begin{itemize}
                \item<2-> \texttt{|} in the column spec: vertical line
                \item<3-> \texttt{\textbackslash hline}: horizontal line
            \end{itemize}
        \end{column}
    \end{columns}
    \note{%
        Show how lines are drawn first, because someone will ask.
        The next slide is why they should mostly not draw them.
    }
\end{frame}

\begin{frame}
    \frametitle{Fewer lines, easier to read}
    \begin{columns}[T, onlytextwidth]
        \begin{column}{.47\textwidth}
            {\large\bfseries\nope{} Grid\par}
            \smallskip
            \tcbinputlst[fontsize=\footnotesize]{listings/04-grid.tex}
            \begin{center}
                \begin{tabular}{|l|l|l|}
    \hline
    Vessel & Material       & Lid \\
    \hline
    Mug    & glazed ceramic & no  \\
    \hline
    Flask  & steel          & yes \\
    \hline
\end{tabular}

            \end{center}
        \end{column}
        \begin{column}{.47\textwidth}
            {\large\bfseries\yep{} \texttt{booktabs}\par}
            \smallskip
            \tcbinputlst[fontsize=\footnotesize]{listings/04-booktabs.tex}
            \begin{center}
                \begin{tabular}{lll}
    \toprule
    Vessel & Material & Lid \\
    \midrule
    Mug & glazed ceramic & no \\
    Flask & steel & yes \\
    \bottomrule
\end{tabular}

            \end{center}
        \end{column}
    \end{columns}
    \note{%
        \begin{itemize}
            \item Point at the two tables and say nothing for three seconds.
                The difference does the teaching.
            \item No vertical lines, no double lines, three horizontal rules of different weight.
                This is how tables in books and journals are set.
            \item Why no vertical lines?
                White space separates columns well enough, and the extra lines make the table harder to scan.
                The \texttt{booktabs} manual explains it in more detail.
        \end{itemize}
    }
\end{frame}

\begin{frame}
    \frametitle{A figure is the same shape}
    \only<1-2,5-| handout:1>{\tcbinputlst[fontsize=\footnotesize]{listings/04-figure.tex}}%
    \only<3-4| handout:none>{\tcbinputlst[fontsize=\footnotesize, highlightlines={3}]{listings/04-figure.tex}}%
    \begin{columns}[T, onlytextwidth]
        \begin{column}{.68\textwidth}
            \begin{itemize}
                \item<2-| handout:1> Same structure as \texttt{table}
                \item<3-| handout:1> \texttt{0.8\textbackslash textwidth}: scales with the page
                \item<4-| handout:1> No file extension
            \end{itemize}
        \end{column}
        \begin{column}{.26\textwidth}
            \includegraphics[width=\textwidth]{renders/cooling-curves.pdf}
        \end{column}
    \end{columns}
    \vfill
    \begin{alertblock}<5->{}
        PDF or PNG, not SVG
    \end{alertblock}
    \note{%
        \begin{itemize}
            \item \texttt{pdflatex} cannot read SVG. Export vector graphics as PDF and they stay sharp.
            \item Without an extension, \LaTeX{} picks the best format it can find.
            \item \texttt{0.8\textbackslash textwidth} survives a change of margins.
            \item ``Same shape as a table'' is the point.
                They learned floats once and get figures nearly free,
                which is why tables come before figures here.
            \item Insist on the no-extension habit:
                it lets the same source build with \texttt{pdflatex} and \texttt{lualatex},
                and it is one less thing to get wrong.
        \end{itemize}
    }
\end{frame}

\begin{frame}
    \frametitle{Pointing at things}
    \tcbinputlst[fontsize=\footnotesize]{listings/04-ref.tex}
    \begin{itemize}
        \item<2-> \texttt{\textasciitilde}: ``Figure'' and its number stay on one line
    \end{itemize}
    \vfill
    \begin{exampleblock}<3->{}
        Move the figure: every number and reference follows
    \end{exampleblock}
    \note{%
        \begin{itemize}
            \item Without the tilde, ``Figure'' can end one line and ``3'' start the next.
                A small thing that marks out people who know what they are doing.
            \item The last line is the payoff of the whole block.
                Pause on it.
        \end{itemize}
    }
\end{frame}

\begin{frame}
    \frametitle{Naming labels}
    \begin{itemize}
        \item Prefixes: \texttt{fig:}, \texttt{tab:}, \texttt{eq:}, \texttt{sec:}
        \item<2-> Sections have labels too:
    \end{itemize}
    \uncover<2->{%
        \tcbinputlst[fontsize=\footnotesize][width=.6\textwidth]{listings/04-section-label.tex}
    }
    \vfill
    \begin{alertblock}<3->{}
        Same line as \texttt{\textbackslash section}
    \end{alertblock}
    \note{%
        \begin{itemize}
            \item Prefixes are only a convention, but they make a long list of labels readable.
            \item A paragraph break between \texttt{\textbackslash section} and \texttt{\textbackslash label}
                can make \texttt{\textbackslash pageref} point at the wrong page.
                Their template has done this since block 1.
        \end{itemize}
    }
\end{frame}

\begin{frame}
    \frametitle{Citations: write the source down once}
    \begin{center}
        \tcbexinputlst[][minted language=bibtex, width=.75\textwidth]{listings/04-bib-entry.bib}[\texttt{references.bib}]
    \end{center}
    \note{%
        \begin{itemize}
            \item Where entries come from matters more than the syntax:
                DOI, Google Scholar, the library's export button, Zotero.
                Nobody types these; say so, half the room assumes they have to.
            \item The key after \texttt{@book\{} is theirs to choose.
        \end{itemize}
    }
\end{frame}

\begin{frame}
    \frametitle{\ldots{} and cite it by name}
    \begin{center}
        \tcbexinputlst[][width=.75\textwidth]{listings/04-bib-use.tex}[In the document]
    \end{center}
    \vfill
    \begin{exampleblock}<2->{}
        \texttt{style=authoryear}: every citation reformats
    \end{exampleblock}
    \note{%
        \begin{itemize}
            \item The style switch is the moment this lands for anyone who ever reformatted a reference list by hand.
            \item \texttt{biblatex} with \texttt{biber}, not legacy Bib\TeX{}.
                A 2008 tutorial will tell them otherwise.
        \end{itemize}
    }
\end{frame}

\begin{frame}
    \frametitle{The one error that still bites}
    \tcbinputlst{listings/04-cite-typo.tex}
    \begin{columns}[T, onlytextwidth]
        \begin{column}{.47\textwidth}
            \uncover<2->{%
                In the PDF:
                \begin{outputbox}
                    \ldots as shown by \textbf{[halliday2019?]}.
                \end{outputbox}
            }
        \end{column}
        \begin{column}{.47\textwidth}
            \uncover<3->{%
                In the log:
                \tcbinputlst[linenos=false][minted language=text, left=2mm, overlay={}]{listings/04-cite-warning.txt}
            }
        \end{column}
    \end{columns}
    \vfill
    \begin{exampleblock}<4->{}
        \texttt{??} or \texttt{[key?]}: the \textbf{name} is wrong
    \end{exampleblock}
    \note{%
        \begin{itemize}
            \item A typo, or a \texttt{\textbackslash label} they never wrote.
                Search for the name in the warning.
            \item TODO 4.5 has them misspell the key deliberately,
                so they meet this in a controlled setting with the log message next to it.
            \item The reflex you are building: read the warning, take the name out of it, search for that name.
                Not ``compile again and hope''.
            \item The other classic cause, too few compilation passes, is handled by Overleaf and \texttt{latexmk}.
                Mention it in one clause if someone asks.
        \end{itemize}
    }
\end{frame}

\begin{frame}
    \frametitle{Your turn}
    \handson{4}{4.1--4.6}{5--10}{checkpoint-04.tex}{%
        The last stretch.\par
        \medskip
        {\normalsize\color{black!60}
        a bare \texttt{tabular} $\rightarrow$ wrap it in a float with \texttt{booktabs}
        $\rightarrow$ place the figure $\rightarrow$ refer to both
        $\rightarrow$ break a citation on purpose, then fix it}
    }
    \note{%
        \textbf{5 to 10 minutes}, and this one genuinely has more in it than that.
        Deliberate: 4.1 to 4.4 is the core, 4.5 and 4.6 are for whoever gets there.
        \begin{itemize}
            \item Say it: ``if you only do the table and the figure, you have got what you came for.''
                It removes the end-of-day panic.
            \item checkpoint-04 is the finished document; anyone who wants the whole thing can take it.
        \end{itemize}
    }
\end{frame}
